% Context from chapters-tex/shannon.tex; for mathematical reference, not compilation.
\section{Continuous and Discrete Signals}

To analyze signal-processing methods, we often begin with a continuous model of the physical input to a sensor.
 
Examples include microphones, digital cameras, and medical imaging devices.
 
For example, a one-dimensional signal can be modeled by $f_0\in\Ldeux([0,1])$, where $[0,1]$ is the acquisition interval, often representing time. A grayscale image can be modeled by $f_0\in\Ldeux([0,1]^2)$ on the unit square.

Although sound and natural images provide many of our examples, most of the methods extend to multidimensional data. After normalizing the signal values, these data can be represented by mappings
\eq{
	f_0 : [0,1]^d \rightarrow [0,1]^s
}
Here $d$ is the dimension of the acquisition domain and $s$ is the number of channels. Grayscale images have $(d,s)=(2,1)$, grayscale videos have $(d,s)=(3,1)$, and RGB color images have $(d,s)=(2,3)$. Multispectral images have $d=2$ and many channels, each recording a different wavelength range. Figures~\ref{fig-examples-1} and~\ref{fig-examples-2} illustrate these types of data.


\myfigure{
	\image{orthobases}{.35}{example-sound}
	\image{orthobases}{.28}{example-image}
	\image{orthobases}{.3}{example-video}
}{ 
	Examples of sound ($d=1$), an image ($d=2$), and a video ($d=3$).  
}{fig-examples-1}

\myfigure{
	\image{orthobases}{.4}{example-color}
	\image{orthobases}{.5}{example-multispectral}
}{ 
	A color image ($s=3$) and a multispectral image ($s=32$).  
}{fig-examples-2}


                                           
\subsection{Acquisition and Sampling}

Signal acquisition maps a continuous signal to a finite-dimensional vector of measurements. A microphone records samples along a time axis, whereas a digital camera records samples on a two-dimensional pixel grid. Abstractly, an acquisition operator maps
\eq{
	f_0 \in \Ldeux([0,1]^d) \mapsto f \in \CC^N
}

\myfigure{
	\image{orthobases}{.4}{discretization-image}
	\image{orthobases}{.5}{discretization-sound}
}{ 
	Image and sound discretization.   
}{fig-discretization}

Figure \ref{fig-discretization} shows examples of discretized signals.

                                           
\subsection{Linear Translation-Invariant Sampling}

To model translation-invariant sampling, extend the signal outside its acquisition domain according to the chosen boundary convention, convolve it with a fixed impulse response $h$, and sample the result:
\eql{\label{eq-linear-sampling}
	f_n = \int_{\RR} f_0(x) h(n/N-x)\,\d x = (f_0\star h)(n/N).
}
The impulse response $h(x)$ depends on the device. It is typically a smooth low-pass kernel concentrated near $x=0$. Its effective width $S$ controls spatial resolution: a wide kernel blurs the signal, whereas a narrow kernel may fail to suppress frequencies that cause aliasing. A width of order $1/N$ balances these effects at the sampling scale.

Section~\ref{subsec-sampling} explains when a bandlimited signal can be reconstructed from its samples. Smoothness alone does not guarantee exact recovery.


                                                                    
                                                                    
                                                                    
\section{Shannon Sampling Theorem}
\label{subsec-sampling}

   
\paragraph{The Fourier transform.}

For $f \in L^1(\RR)$, its Fourier transform is defined as
\eql{\label{eq-fourier-transform}
	\foralls \om \in \RR, \quad
	\hat f(\om) \eqdef \int_\RR f(x) e^{-\imath x \om} \d x.
}
For $f\in L^1(\RR)\cap L^2(\RR)$, Plancherel's identity gives $\norm{\hat f}_2^2 = 2\pi\norm{f}_2^2$. The map $f \mapsto \hat f$ therefore extends continuously to $L^2(\RR)$. This extension is the $L^2$ limit, as $T \rightarrow +\infty$, of the truncated integrals $\int_{-T}^T f(x) e^{-\imath x \om} \d x$.
 
When $\hat f \in L^1(\RR)$, Fourier inversion gives
\eql{\label{eq-i-ft}
	f(x) = \frac{1}{2\pi} \int_\RR \hat f(\om) e^{\imath x \om} \d \om, 
}
where the equality holds almost everywhere. The right-hand side defines the con